%=============================================================================
%  CHAPTER 6 — IMPLEMENTATION
%=============================================================================
\chapter{Implementation}

\section*{Chapter Introduction}

This chapter presents the implementation of the ForsaDrive system based on the design defined in the previous chapter. It describes the development environment, project organisation, and the implementation of core system components, including the backend, web application, and mobile application. Additionally, it highlights the integration between these components and discusses the main technical challenges encountered during development along with the solutions adopted.

%─────────────────────────────────────────────────────────────────────────────
\section{Development Environment}

The development of ForsaDrive was carried out using a set of tools and technologies selected to ensure productivity, flexibility, and compatibility across platforms.

\subsection{Hardware Environment}

\begin{itemize}[noitemsep]
    \item Development machines: Personal computers (Windows/Linux)
    \item Minimum specifications: 8GB RAM, multi-core CPU
\end{itemize}

\subsection{Software Environment}

\begin{itemize}[noitemsep]
    \item Code Editor: Visual Studio Code
    \item Backend Runtime: PHP
    \item Mobile Development: Flutter SDK
    \item Database: SQLite
    \item Version Control: Git
    \item API Testing: Postman
\end{itemize}

These tools enabled efficient development, testing, and debugging of both web and mobile components.

%─────────────────────────────────────────────────────────────────────────────
\section{Project Structure}

The project is organised into three main components:

\begin{itemize}[noitemsep]
    \item \textbf{Backend:} Handles business logic and database interaction
    \item \textbf{Web Application:} Provides a full-featured interface
    \item \textbf{Mobile Application:} Offers a mobile-friendly user experience
\end{itemize}

\subsection{Backend Structure}

The backend follows a modular structure inspired by MVC principles:

\begin{itemize}[noitemsep]
    \item \textbf{Controllers:} Handle incoming requests
    \item \textbf{Models:} Manage database operations
    \item \textbf{Views:} Render web pages
\end{itemize}

\subsection{Mobile Structure}

The Flutter project is organised into:

\begin{itemize}[noitemsep]
    \item Screens (UI pages)
    \item Providers (state management)
    \item Services (API communication)
    \item Models (data representation)
\end{itemize}

This separation improves maintainability and readability.

%─────────────────────────────────────────────────────────────────────────────
\section{Backend Implementation}

The backend serves as the core of the system, implementing all business logic and exposing RESTful APIs.

\subsection{Authentication System}

The authentication module manages:
\begin{itemize}[noitemsep]
    \item User registration
    \item Login and session management
    \item Password hashing and validation
\end{itemize}

Security is ensured through input validation and protected access to sensitive endpoints.

\subsection{Ride Management}

Drivers can publish rides by specifying:
\begin{itemize}[noitemsep]
    \item origin and destination
    \item date and time
    \item number of available seats
    \item price per seat
\end{itemize}

The system validates all inputs before storing ride data.

\subsection{Booking System}

The booking module handles:
\begin{itemize}[noitemsep]
    \item seat reservation
    \item availability checks
    \item wallet balance validation
\end{itemize}

Bookings are only confirmed if all conditions are satisfied.

\subsection{Messaging System}

Users can communicate through an integrated messaging system:
\begin{itemize}[noitemsep]
    \item conversation threads between users
    \item message storage and retrieval
\end{itemize}

\subsection{Payment and Wallet System}

The wallet system manages:
\begin{itemize}[noitemsep]
    \item balance updates
    \item booking deductions
    \item refunds
    \item referral rewards
\end{itemize}

%─────────────────────────────────────────────────────────────────────────────
\section{Web Application Implementation}

The web application provides a complete interface for both users and administrators.

\subsection{User Interface}

The interface is designed using Bootstrap, ensuring responsiveness and consistency.

\subsection{Key Functional Pages}

\begin{itemize}[noitemsep]
    \item Authentication pages (login/register)
    \item Dashboard
    \item Ride search and booking
    \item Ride publishing
    \item Messaging interface
    \item Admin panel
\end{itemize}

\begin{figure}[H]
\centering
\includegraphics[width=\textwidth]{web_admin.png}
\caption{Web Administration Panel}
\end{figure}

The web application also includes administrative tools for managing users, rides, complaints, student verifications, driver applications, organizations, and student domain lists across seven dedicated tabs.

%─────────────────────────────────────────────────────────────────────────────
\section{Mobile Application Implementation}

The mobile application extends the platform to smartphone users.

\subsection{Core Features}

\begin{itemize}[noitemsep]
    \item User authentication
    \item Ride search and booking
    \item Ride management
    \item Messaging
    \item Profile management
\end{itemize}

\subsection{State Management}

The application uses the Provider pattern to manage state and ensure efficient UI updates.

\subsection{API Integration}

All data operations are performed through REST API calls to the backend using the Dio package.

\begin{figure}[H]
\centering
\includegraphics[width=0.6\textwidth]{mobile_home.png}
\caption{Mobile Application Home Screen}
\end{figure}

%─────────────────────────────────────────────────────────────────────────────
\section{Integration Between Components}

The web and mobile applications both interact with the same backend through REST APIs.

\begin{itemize}[noitemsep]
    \item Shared business logic ensures consistency
    \item Centralised database guarantees data integrity
    \item API endpoints standardise communication
\end{itemize}

This architecture enables seamless interaction across platforms.

%─────────────────────────────────────────────────────────────────────────────
\section{Challenges and Solutions}

During development, several technical challenges were encountered:

\begin{itemize}[noitemsep]
    \item \textbf{API Synchronisation:} Ensuring consistent behaviour between web and mobile  
    \textit{Solution: Unified API design}
    
    \item \textbf{State Management (Mobile):} Managing UI updates efficiently  
    \textit{Solution: Provider architecture}
    
    \item \textbf{Security:} Preventing invalid inputs and unauthorised access  
    \textit{Solution: Input validation and role-based control}
    
    \item \textbf{Database Limitations:} SQLite constraints  
    \textit{Solution: Structured schema design and optimisation}
\end{itemize}

These solutions improved system stability and reliability.

%─────────────────────────────────────────────────────────────────────────────
\section*{Chapter Conclusion}

This chapter detailed the implementation of the ForsaDrive system, covering backend development, web and mobile applications, and their integration. It also highlighted the main challenges encountered and the solutions adopted.

The following chapter presents the testing, validation, and deployment of the system.